\chapter{Implementation}

Mit der beschriebenen Konzeption kann nun die Implementation des Webapplikation
besprochen werden. Diese wird dabei in das Frontend und Backend aufgeteilt. So
wurde zwar immer Feature basiert gearbeitet, \dh{} um eine Funktionalität zu
erweitern mussten meistens sowohl das Frontend als auch das Backend angepasst
werden, jedoch hat man eine starke Konzeptionelle Unterscheidung, wie beide
Applikationsteile aufgebaut sind. So sollen in den folgenden Abschnitten
zuerst das Backend und danach das Frontend beschrieben werden. Zuletzt soll noch eine
kurze Übersicht gegeben werden, wie die Webapplikation ausgeführt werden kann.

\section{Backend}

Das Backend wurde, wie schon in der Konzeption erläutert, mithilfe von der
Programmiersprache \emph{Kotlin} geschrieben und folgt einer
Schichtenarchitektur mit insgesamt drei Ebenen. Diese Schichten werden in den
folgendens Unterkapiteln näher betrachtet.

\subsection{Datenschicht}

Die Datenschicht soll es ermöglichen, einen Zugriff auf Daten von verschiedenen
Persistenz-Technologien zu erhalten und diese zu manipulieren. Dafür wurde hauptsächlich
mit der beschriebenen relationalen Datenbank PostgreSQL und der Bibliothek \emph{sqldelight}
gearbeitet. Durch diese Bibliothek konnte das Datenmodell und die Queries durch
SQL definiert werden und die Bibliothek generierte daraus Typsicheren Quelltext zur Interaktion
mit der Datenbank. Dafür müssen zuerst Migrationsdateien, welche einfache SQL-Skripte sind,
geschrieben werden. In diesen werden die üblichen Statements wie \code{CREATE TABLE}
das Datenmodell definiert. Daraus kann \emph{sqldelight} aus den definierten Tabellen einfache
Datenklassen zu den bilden. Separat werden auch noch die Queries für die Interaktion mit den Tabellen
definiert. Dafür werden Dateien erstellt aus denen \code{Query} Klassen generiert werden. In dieser
können Queries geschrieben werden, zu welchen immer ein Label gesetzt wird, welche den späteren
Methodennamen darstellt.

Mit der Quelltext-Generation konnte dann der \code{DatabaseService} geschrieben werden. In diesem
wird eine Verbindung über einen Connection-Pool zur Datenbank herstellt und mit der bestehenden Verbindung
die Query-Klassen initialisiert. Zusätzlich besteht ein \code{MigrationService}, der im \code{DatabaseService}
genutzt wird. So müssten die bereits in SQL geschriebenen Migrationen auch ausgeführt werden, und die bisherige
Version des Datenbankschemas persistiert werden. \emph{sqldelight} generiert dafür eine \code{migrate}-Methode,
die die Migrationen ausführt. Für diese muss aber die bisherige Version des Datenbankschemas übergeben werden.
Dafür wird im \code{MigrationService} eine Tabelle in der Datenbankerstellt, welche die Version persistiert.
So kann diese Version für die Migration erhalten werden und nach der Migration auf die neueste Version gesetzt werden.

Neben der Datenbank wird noch im Backend der geteilter In-Memory-Cache \emph{redis}~\cite{redis} eingesetzt.
Damit können Daten geteilt bei mehreren Backend-Instanzen zwischengespeichert werden. Dies wurde verwendet um
Datenbankabfragen, welche eine längere Zeit dauern zwischenzuspeichern und die Ergebnisse schneller zurückgeben
zu können. Ein Beispiel ist hierfür der \code{AuthorizationCache}, in welcher die Zugriffsberechtigungen für
das Absichern der API zwischenspeichert.

\subsection{Businessschicht}

Mit dem rohen Datenzugriff kann nun die eigentliche Applikationslogik für das
Backend erstellt werden. Dies erfolgt in den \emph{Services} Klassen, welche
auf den \emph{Queries} und \emph{Caches} der vorherigen Schicht aufbauen. Diese enthalten die
spezifischen Applikationslogik, um \zb{} die Aktionen für die Teams oder der
Authentifikation durchzuführen. Diese Services haben dabei keinen
internen Zustand und hängen nur von der Datenhaltung oder weiteren Services ab.

Die Services für den Nutzer, die Teams, die Team-Kanäle und die Team-Teilnehmer, bieten dabei
eine sehr ähnliche Schnittstelle an. Diese hauptsächlich ein Create-Read-Update-Delete (CRUD)
Interface an, um eines oder mehrere Datensätze zu erhalten, neue zu erstellen
oder alte zu verändern und zu löschen. Es gibt auch spezifischere Aktionen auf den jeweiligen
Daten, welches \zb{} welches mehr als nur die Manipulation eines Objektes
erfordert. Dafür werden domänenspezifische Aktionen angeboten. Ein Beispiel ist das Beitreten
eines neuen Teammitgliedes zu einem Team. Hier werden auch die Daten aus der Persistenz-Schicht
zu domänenspezifischen Objekten umgewandelt, welche auch für die spätere Präsentationsschicht
serialisierbar sind.

So gibt es aber auch Services, die nicht in das CRUD-Schema passen. Dies ist
\zb{} der \code{AuthenticationService} oder \code{AuthorizationService}, welche die Nutzer Authentifikation und
Autorisation steuert. Über den \code{Authentication\-Service} können Nutzerkonten
über eine Registration erstellt werden. Um den Nutzer für die spätere Authentifikation zu
identifizieren wird eine Kombination aus einem kurzlebigen Access-Token und
einem langlebigen Refresh-Token verwendet. Der Access-Token, ist dabei ein
signierter \emph{JWT} Token, welcher nur eine kurze Gültigkeit von einem Tag
hat. Dieser wird zu der eigentlichen Authentifikation verwendet. Wenn der
bisherige Access-Token abgelaufen ist, kann der Refresh-Token verwendet werden,
um einen neuen AccessToken generieren zu lassen. Dafür muss der Refresh-Token
in der Datenbank gespeichert werden. Durch diese Kombination beider Tokens, muss
nur selten ein Datenbankzugriff zur Authentifikation erfolgen. Zuletzt hat der
\code{AuthorizationService} Methoden zur Autorisation. Mit diesen kann geprüft werden,
ob ein Nutzer der Besitzer eines Projektes ist oder dieses für ihn geteilt
wurde, womit in der Präsentationsschicht die API-Routen abgesichert werden können.

Das das Backend über eine Echtzeit-Komponente für Chats der einzelnen Teams und zur
generellen Benachrichtigung von Events verfügt, gibt es noch \emph{Services} zur Steuerung
von diesen. Um auch hier eine Skalierung des Backends zu ermöglichen, wurde das Pub-Sub-System \emph{Nats}~\cite{nats}
eingesetzt. Mit diesem System ist es möglich sich für verschiedene Nachrichtenkanäle zu registrieren,
und beim Senden von Nachrichten davon benachrichtigt werden. So kann \zb{} bei mehreren
Backend-Instanzen eine Instanz für den Client eine Chat-Nachricht senden und alle weiteren
Backends würden davon benachrichtigt werden. Das wird verwendet, da die Websocket-Verbindungen
zwar zustandsvoll sind und so immer an ein Backend gebunden sind, jedoch mit dieser Technik
Benachrichtigungen, immer an alle Clients ausgeliefert werden können, egal mit welcher spezifischen
Instanz, diese sich verbunden haben.

\subsection{Präsentationsschicht}

Als letzte Schicht wird noch die Präsentationsschicht betrachtet. Diese bietet
eine REST-Schnittelle durch die Bibliothek \emph{ktor} an. Zum Erstellen von
API-Routen wurden Klassen erstellt, welche das selbst geschrieben \code{HttpRouter}
Interface implementieren. In diesem wird ein Callback erfüllt, über welchen die einzelnen
API-Routen registriert werden können. Für jede Route wird so in der Klasse eine
Handler-Methode geschrieben. In dieser Methode stehen einen die verschiedenen
HTTP-Request Daten zur Verfügung. So kann in den Methoden \zb{} ein JSON-Objekt aus dem
HTTP-Body der Abfrage deserialisiert werden. Mit dem erhaltenen Daten können dann
die Methode der Services aufgerufen, die erhaltenen domänenspezifischen Objekte
in JSON-Objekte serialisiert und diese darauf als Antwort dem Client zurückgegeben werden.
Die implementierten \code{HttpRouter} konnten dann dem \code{HttpService} übergeben werden,
wo der Http-Server von \code{ktor} gestartet wurde um den Router-Klassen konfiguriert  wurde.

\subsection{Dependecy Injection}

Für eine automatische Instanziierung aller benötigter Klassen und eine passende
Anbindung derer Abhängigkeiten, wurde die Dependency-Injection Bibliotheken
\emph{dagger} und \emph{anvil} verwendet. Für die Integration in die bisher beschriebenen
Klassen müssen hierfür nur um Annotationen ergänzt werden. Dadurch wird mitgeteilt, über welchen
Konstruktor \zb{} die Klasse initialisierbar ist.

Zum Erstellen und Initialisieren des gesamten Abhängigkeits-Graphen der Applikation stellt der \code{AppComponent}
den Einttritspunkt dar. Für dieses erstellt \emph{dagger} den gesammten Injektions-Quelltext, um
eine konkrete Instanz des \code{AppComponents} zu erhalten. In diesem ist der \code{ServiceManager}
vorhanden, welcher die Dienste der Applikation verwaltet und startet.

\section{Frontend}

Nun sollte auch das Frontend der Applikation beschrieben werden, welches mit
der Programmiersprache \emph{TypeScript} und der UI-Bibliothek \emph{SolidJS}
entwickelt wurde. Das Frontend hat dabei die Aufgabe, Daten von dem Backend zu
erhalten und daraus eine Interaktive-Nutzeroberfläche im Browser zu erstellen.
Auch sollten bei Nutzerinteraktionen die Daten des Backends manipuliert werden
können. Dafür kann das Frontend in zwei Teile gegliedert werden: Die
UI-Komponenten, welche die UI und deren Interaktivität bilden und den
dazugehörigen Zuständen, welche die Daten zur Darstellung und Methoden für die
Datenmanipulation für die UI-Komponenten anbieten.

Da die einzelnen Zustände die Grundlage für die UI-Komponenten bilden, sollten
sich diese zuerst betrachtet werden. Da \emph{SolidJS} eine reaktive
UI-Bibliothek ist, werden für das Bilden von Zuständen, primitive
Datenstrukturen angeboten, mit welchen die UI-Komponenten angepasst werden
können, falls sich der Inhalt der Datenstrukturen ändern. Ein Beispiel ist
dafür die \code{createSignal} Funktion. Die einen Getter für den Erhalt der
Daten und einen Setter zum Setzen neuer Daten zurückgibt. Wenn ein neuer Wert
über den Setter gesetzt wird, wird jede Komponente, die den Getter verwendet
mit dem neuen Inhalt angepasst. Solche primitive werden nun in
\emph{State}-Klassen verwendet, um den Zustand für die UI zu bilden. Dafür
halten diese zum einen den Zustand für die UI-Logik, wie \zb{} den Zustand
einer Form für das Erstellen das Verwalten eines Teams. Auch halten die States, die aus dem
Backend erhaltenen Daten. Die Zustandsklasse bietet nun eine Schnittstelle
durch Methoden an, um die internen Daten zu erhalten und zu verändern.

Diese Zustände werden verwendet, um die UI zu bilden und eine Interaktivität zu
dieser hinzuzufügen. Dafür werden in \emph{SolidJS} \sog{} Komponenten
verwendet, welche einfache Funktionen sind, die HTML-Markup oder andere
Komponenten zurückgeben können. Diese werden im Vergleich zu anderen
Komponentenbasierten UI-Bibliotheken nur einmal für das Rendern aufgerufen. So
kann hierarchisch die gesamte Applikation gebildet werden. Die Komponenten sind
dabei nach den möglichen Navigations-Routen der Applikation organisiert. Diese
können in der \code{App} Komponenten zu der Applikation über einen Routing
Mechanismus eingebunden werden. So gibt es immer für jede mögliche Destination
eine \code{Route} Komponente, die über diese Route verantwortlich ist. Dabei
können auch Routen weitere Routen selbst erhalten. So ist \zb{} die Route
für ein einzelnes Team in der Route für die gesamte Applikation eingebunden.
So werden wie die Komponenten selbst, die Routen hierarchisch angeordnet In
diesen Routen werden auch immer die passenden Zustände aus dem vorherigen Abschnitt
gebildet. Alle weiteren Komponenten unter einer \code{Route} können dabei den Zustand
von der Route erhalten, um mit den Daten die UI zu bilden oder die Interaktivität
der UI zu geben.

\section{Inbetriebnahme}

Zuletzt sollte noch beschrieben werden, wie das Projekt gestartet werden kann.
Hierfür können werden zwei Möglichkeiten geboten. So kann man hierfür zum einen
die lokale Entwicklungsvariante verwenden. In dieser müssen das Frontend und
Backend separat gestartet werden. Eine genaue Anleitung findet sich in der
\code{README} Datei im Root des Projektes. Hierbei kann mit der Entwicklungsumgebung
\emph{Intellij IDEA} das Backend gestartet werden. Das Frontend hingegen benötigt,
dadurch dass es in \emph{TypeScript} geschrieben wurde, eine relativ neue Version von
\emph{Node.js}~\cite{node} und der Paket-Manager \emph{pnpm}~\cite{pnpm}.
Diese Variante kann verwendet werden, um das Projekt auszuführen jedoch sollte dafür
wegen der Einfachheit die nächste Variante verwendet werden.

Eine weitere Möglichkeit, das Projekt in Betrieb zu nehmen, ohne die gesamten
Abhängigkeiten für das Frontend und Backend zu installieren, ist das Container
Orchestrierungstool \emph{Docker}~\cite{docker}. Dafür wird nur \emph{docker}
selbst und das Tool \emph{docker-compose} benötigt, welches bei den älteren
Versionen von \emph{docker} seperat installiert werden muss. Dafür findet sich
im Root des Projektes eine \code{compose.yml} Datei. In dieser wird das
Backend und Frontend neben den Persistenz-Technologien gestartet. Hierbei sollten
noch Konfigurationsdateien erstellt werden. Dafür braucht es im Root des Projektes
eine \code{.env} und eine \code{config.json} Datei. Auch hierfür findet sich in der \code{README}
Datei des Projektes eine genaue Anleitung, wie mit welchen Daten diese Dateien befüllt werden sollten.
Danach kann das Projekt mit den \code{docker-compose up ----build} Befehl gestartet werden.
